\subsection{DNS 2020 \cite{reddy2020}}
The Deep Noise Suppression challenge is a yearly challenge organized by Microsoft Research. The first challenge was introduced in 2020, and it has been held annually since. If focuses specially on Speech Enhancement tasks and real-time settings. The Deep Noise Suppression (DNS) challenge aims to unify research activity in the Speech Enhancement domain by making the train/test datasets and subjective assessment system available to the public. It releases a large clean speech and noise datasets. It also provides configurable synthetization scripts for producing the training datasets. It also releases a part of the test dataset.

\subsubsection{Clean Dataset}
The clean speech dataset is obtained from Librivox, which a public audiobooks dataset. It has the audiobooks recordings of over 10000 public domain books in multiple languages, the majority of which are in English \cite{panayotov2015}. Although most of these recordings were of excellent speech conditions and quality, it had many recordings with poor conditions such as speech distortion, background noise, and reverberation. These poor-conditioned recordings needed to be removed, as the clean dataset must be free of noisy and poor-conditioned recordings. An online subjective test framework ITU-T P.808 was used to keep only excellent quality recordings. Mean Opinion Score (MOS) was used as a metric by keeping only the upper quartile with respect to MOS. The final clean dataset included 500 hours of speech from 2150 speakers. The filtered clips are then divided into 10-second segments.

\subsubsection{Noise Dataset}
The noise dataset was selected from Audioset \cite{gemmeke2017} and Freesound. Audioset consists of 2 million sound clips of 10s each, where each sound clip is human-labelled. These clips are extracted from YouTube videos belonging to about 600 audio events. Some classes are overrepresented compared to others. A sampling approach was executed to balance each class by having at least 500 clips. Speech activity detector was used to remove any clips containing speech activity. The resulting dataset contained about 150 classes and 60000 clips. Additional 10000 noise clips related to VoIP applications were downloaded from Freesound and DEMAND dataset \cite{thiemann2013}.

\subsubsection{Noisy Dataset}
The noisy dataset is the addition of clean speech from the clean dataset and a noise clip from the noise dataset at various SNR levels. The SNR is only computed using active segments of both the clean and noise signals. This was done to avoid the overshoot of amplitude in impulsive noise (i.e., door shutting or dog barking).
